% Фрагмент магистерской диссертации: цель, задачи и структура работы.
% Тема: «Создание имитационной модели для анализа поведения ансамбля
% беспилотных мобильных платформ. Развёртка имитационной модели с
% масштабируемыми параметрами».
%
% Подключение к основному документу: % Фрагмент магистерской диссертации: цель, задачи и структура работы.
% Тема: «Создание имитационной модели для анализа поведения ансамбля
% беспилотных мобильных платформ. Развёртка имитационной модели с
% масштабируемыми параметрами».
%
% Подключение к основному документу: % Фрагмент магистерской диссертации: цель, задачи и структура работы.
% Тема: «Создание имитационной модели для анализа поведения ансамбля
% беспилотных мобильных платформ. Развёртка имитационной модели с
% масштабируемыми параметрами».
%
% Подключение к основному документу: % Фрагмент магистерской диссертации: цель, задачи и структура работы.
% Тема: «Создание имитационной модели для анализа поведения ансамбля
% беспилотных мобильных платформ. Развёртка имитационной модели с
% масштабируемыми параметрами».
%
% Подключение к основному документу: \input{отчётность/magister_tsel_zadachi_struktura}
% либо автономная сборка: pdflatex magister_tsel_zadachi_struktura.tex
%
\documentclass[12pt,a4paper]{article}
\usepackage[utf8]{inputenc}
\usepackage[T2A]{fontenc}
\usepackage[russian]{babel}
\usepackage{geometry}
\geometry{left=30mm,right=15mm,top=20mm,bottom=20mm}
\usepackage{indentfirst}
\setlength{\parindent}{1.25cm}

\begin{document}

\section*{Цель исследования}

\textbf{Целью} магистерской диссертации является разработка и развёртка
имитационной модели ансамбля беспилотных мобильных платформ,
позволяющей воспроизводить типовые режимы группового полёта и
анализировать устойчивость, координацию и качество управления при
изменении числа агентов, геометрии постановки задачи и параметров
законов управления и телеметрии (в том числе в условиях ограничений
канала связи и шумов измерений), с оформлением воспроизводимого
конвейера экспериментов и представлением результатов в виде
логов, метаданных прогона и графической визуализации.

\section*{Задачи исследования}

Для достижения указанной цели в работе ставятся следующие \textbf{задачи}:
\begin{enumerate}
  \item Провести обзор литературы и нормативно\-технической базы по
    моделированию многоагентных БПЛА, программно\-интегрированным
    стендам SITL/HIL, протоколу MAVLink и методам координации роя
    (лидер\,--\,последователь, формации, локальные законы соседства).
  \item Сформулировать постановку задачи имитации ансамбля: состав
    подсистем (динамика носителя в среде симуляции, обмен данными,
    сценарий верхнего уровня), требования к масштабируемости
    (число платформ, сетевые порты, начальные положения, длительность
    эксперимента, файлы параметров автопилота).
  \item Обосновать и описать используемые методы: представление
    состояния в локальной системе координат (NED), извлечение
    оценок положения из телеметрии (включая SIM\_STATE и связанные
    сообщения), преобразования «географические координаты\,--\,NED»,
    организация потокобезопасного доступа к соединению MAVLink,
    алгоритмы законов управления в сценариях (в частности,
    пропорционально\-интегральные и кусочно\-непрерывные/сигмоидные
    отображения ошибки в команды RC override), критерии и метрики
    анализа (расстояния между агентами, скорости, PWM, параметры
    устойчивости относительно опорных узлов).
  \item Реализовать программный комплекс имитации: запуск нескольких
    экземпляров SITL (ArduPilot), сценарии на языке Python,
    ядро взаимодействия с автопилотом по MAVLink, режимы с
    опциональной 3D\-визуализацией (Webots) и 2D\-визуализацией,
    пакетный запуск серии экспериментов с фиксацией метаданных.
  \item Выполнить серию вычислительных экспериментов по выбранным
    сценариям (включая сценарии сравнения при наличии шума в канале
    или в оценке положения), верифицировать корректность развёртки
    и воспроизводимость результатов.
  \item Систематизировать полученные результаты, оформить выводы
    о применимости разработанной имитационной модели для анализа
    поведения ансамбля и направлениях дальнейшего развития
    (расширение класса сценариев, учёт более сложных моделей среды
    и связи).
\end{enumerate}

\section*{Структура работы}

Диссертация состоит из введения, трёх глав, заключения, списка
использованных источников и приложений. Ниже приведена логическая
связка глав: в главе~2 излагаются методы и модели, в главе~3\,---
их программная реализация и результаты экспериментов по тем же
разделам (теория\,--\,практика).

\textbf{Во введении} обосновывается актуальность темы, формулируются
цель и задачи, кратко характеризуется объект и предмет исследования,
перечисляются научная новизна и практическая значимость, приводится
обзор структуры работы и положений, выносимых на защиту.

\textbf{Глава~1. Обзор литературы в рамках диссертации}
\begin{enumerate}
  \item Состояние исследований по координации ансамблей БПЛА:
    централизованные и децентрализованные подходы, устойчивость
    и ограничения по связи.
  \item Программно\-аппаратные стенды и имитация: SITL/HIL,
    интеграция с Gazebo/Webots, типовые стеки автопилота.
  \item Протокол MAVLink, модели телеметрии и команд верхнего
    уровня; вопросы синхронизации и частоты обновления данных.
  \item Методы анализа результатов имитации: метрики близости
    к заданной конфигурации, безопасные интервалы, сравнение
    вариантов алгоритма и параметров.
\end{enumerate}

\textbf{Глава~2. Обзор использованных методов, математики и алгоритмов}
\begin{enumerate}
  \item Математическое описание кинематики и используемых систем
    координат; переход от географических координат к локальному NED.
  \item Модель информационного обмена в имитационном стенде:
    потоки MAVLink, выбор сообщений для оценки состояния,
    организация параллельного доступа (потокобезопасный обработчик
    соединения).
  \item Класс рассматриваемых законов управления ансамблем:
    сценарии лидер\,--\,последователь, формации, цепочки с опорными
    узлами, локальные законы «антенна/соседство»; параметризация
    законов (коэффициенты ПИ, ограничения PWM, радиусы видимости,
    шаг дискретизации по PWM и~т.\,п.).
  \item Постановка вычислительного эксперимента: воспроизводимость,
    масштабирование числа агентов, пакетные прогоны, фиксация
    метаданных и логов для последующего анализа.
  \item Методы оценки качества: показатели по траекториям, отклонениям
    от якоря, согласованности PWM, сравнение режимов «без шума /
    со шумом» (при наличии соответствующего сценария).
\end{enumerate}

\textbf{Глава~3. Практическая реализация алгоритмов, фрагменты кода,
визуализация и результаты}
\begin{enumerate}
  \item Архитектура программного комплекса: сценарный уровень,
    модуль взаимодействия с MAVLink, запуск SITL и распределение
    сетевых портов; конфигурация параметров автопилота.
  \item Реализация геометрических и навигационных преобразований
    и потокобезопасного приёма/передачи сообщений (иллюстрация
    ключевыми фрагментами кода и пояснениями).
  \item Реализация выбранных сценариев ансамбля: соответствие
    алгоритмам главы~2, настраиваемые параметры сценария,
    журналирование телеметрии в файлы эксперимента.
  \item Средства визуализации (2D/3D) и конвейер постобработки:
    построение графиков траекторий, скоростей, PWM и других
    величин по сохранённым логам.
  \item Результаты вычислительных экспериментов и их интерпретация
    в контексте цели работы; обсуждение ограничений модели.
\end{enumerate}

\textbf{В заключении} формулируются основные результаты, степень
достижения цели, практическая значимость и перспективы развития.

\textbf{В приложениях} могут быть приведены листинги конфигураций,
таблицы параметров прогонов (включая JSON\-метаданные экспериментов),
дополнительные графики и описание сценариев запуска.

\end{document}

% либо автономная сборка: pdflatex magister_tsel_zadachi_struktura.tex
%
\documentclass[12pt,a4paper]{article}
\usepackage[utf8]{inputenc}
\usepackage[T2A]{fontenc}
\usepackage[russian]{babel}
\usepackage{geometry}
\geometry{left=30mm,right=15mm,top=20mm,bottom=20mm}
\usepackage{indentfirst}
\setlength{\parindent}{1.25cm}

\begin{document}

\section*{Цель исследования}

\textbf{Целью} магистерской диссертации является разработка и развёртка
имитационной модели ансамбля беспилотных мобильных платформ,
позволяющей воспроизводить типовые режимы группового полёта и
анализировать устойчивость, координацию и качество управления при
изменении числа агентов, геометрии постановки задачи и параметров
законов управления и телеметрии (в том числе в условиях ограничений
канала связи и шумов измерений), с оформлением воспроизводимого
конвейера экспериментов и представлением результатов в виде
логов, метаданных прогона и графической визуализации.

\section*{Задачи исследования}

Для достижения указанной цели в работе ставятся следующие \textbf{задачи}:
\begin{enumerate}
  \item Провести обзор литературы и нормативно\-технической базы по
    моделированию многоагентных БПЛА, программно\-интегрированным
    стендам SITL/HIL, протоколу MAVLink и методам координации роя
    (лидер\,--\,последователь, формации, локальные законы соседства).
  \item Сформулировать постановку задачи имитации ансамбля: состав
    подсистем (динамика носителя в среде симуляции, обмен данными,
    сценарий верхнего уровня), требования к масштабируемости
    (число платформ, сетевые порты, начальные положения, длительность
    эксперимента, файлы параметров автопилота).
  \item Обосновать и описать используемые методы: представление
    состояния в локальной системе координат (NED), извлечение
    оценок положения из телеметрии (включая SIM\_STATE и связанные
    сообщения), преобразования «географические координаты\,--\,NED»,
    организация потокобезопасного доступа к соединению MAVLink,
    алгоритмы законов управления в сценариях (в частности,
    пропорционально\-интегральные и кусочно\-непрерывные/сигмоидные
    отображения ошибки в команды RC override), критерии и метрики
    анализа (расстояния между агентами, скорости, PWM, параметры
    устойчивости относительно опорных узлов).
  \item Реализовать программный комплекс имитации: запуск нескольких
    экземпляров SITL (ArduPilot), сценарии на языке Python,
    ядро взаимодействия с автопилотом по MAVLink, режимы с
    опциональной 3D\-визуализацией (Webots) и 2D\-визуализацией,
    пакетный запуск серии экспериментов с фиксацией метаданных.
  \item Выполнить серию вычислительных экспериментов по выбранным
    сценариям (включая сценарии сравнения при наличии шума в канале
    или в оценке положения), верифицировать корректность развёртки
    и воспроизводимость результатов.
  \item Систематизировать полученные результаты, оформить выводы
    о применимости разработанной имитационной модели для анализа
    поведения ансамбля и направлениях дальнейшего развития
    (расширение класса сценариев, учёт более сложных моделей среды
    и связи).
\end{enumerate}

\section*{Структура работы}

Диссертация состоит из введения, трёх глав, заключения, списка
использованных источников и приложений. Ниже приведена логическая
связка глав: в главе~2 излагаются методы и модели, в главе~3\,---
их программная реализация и результаты экспериментов по тем же
разделам (теория\,--\,практика).

\textbf{Во введении} обосновывается актуальность темы, формулируются
цель и задачи, кратко характеризуется объект и предмет исследования,
перечисляются научная новизна и практическая значимость, приводится
обзор структуры работы и положений, выносимых на защиту.

\textbf{Глава~1. Обзор литературы в рамках диссертации}
\begin{enumerate}
  \item Состояние исследований по координации ансамблей БПЛА:
    централизованные и децентрализованные подходы, устойчивость
    и ограничения по связи.
  \item Программно\-аппаратные стенды и имитация: SITL/HIL,
    интеграция с Gazebo/Webots, типовые стеки автопилота.
  \item Протокол MAVLink, модели телеметрии и команд верхнего
    уровня; вопросы синхронизации и частоты обновления данных.
  \item Методы анализа результатов имитации: метрики близости
    к заданной конфигурации, безопасные интервалы, сравнение
    вариантов алгоритма и параметров.
\end{enumerate}

\textbf{Глава~2. Обзор использованных методов, математики и алгоритмов}
\begin{enumerate}
  \item Математическое описание кинематики и используемых систем
    координат; переход от географических координат к локальному NED.
  \item Модель информационного обмена в имитационном стенде:
    потоки MAVLink, выбор сообщений для оценки состояния,
    организация параллельного доступа (потокобезопасный обработчик
    соединения).
  \item Класс рассматриваемых законов управления ансамблем:
    сценарии лидер\,--\,последователь, формации, цепочки с опорными
    узлами, локальные законы «антенна/соседство»; параметризация
    законов (коэффициенты ПИ, ограничения PWM, радиусы видимости,
    шаг дискретизации по PWM и~т.\,п.).
  \item Постановка вычислительного эксперимента: воспроизводимость,
    масштабирование числа агентов, пакетные прогоны, фиксация
    метаданных и логов для последующего анализа.
  \item Методы оценки качества: показатели по траекториям, отклонениям
    от якоря, согласованности PWM, сравнение режимов «без шума /
    со шумом» (при наличии соответствующего сценария).
\end{enumerate}

\textbf{Глава~3. Практическая реализация алгоритмов, фрагменты кода,
визуализация и результаты}
\begin{enumerate}
  \item Архитектура программного комплекса: сценарный уровень,
    модуль взаимодействия с MAVLink, запуск SITL и распределение
    сетевых портов; конфигурация параметров автопилота.
  \item Реализация геометрических и навигационных преобразований
    и потокобезопасного приёма/передачи сообщений (иллюстрация
    ключевыми фрагментами кода и пояснениями).
  \item Реализация выбранных сценариев ансамбля: соответствие
    алгоритмам главы~2, настраиваемые параметры сценария,
    журналирование телеметрии в файлы эксперимента.
  \item Средства визуализации (2D/3D) и конвейер постобработки:
    построение графиков траекторий, скоростей, PWM и других
    величин по сохранённым логам.
  \item Результаты вычислительных экспериментов и их интерпретация
    в контексте цели работы; обсуждение ограничений модели.
\end{enumerate}

\textbf{В заключении} формулируются основные результаты, степень
достижения цели, практическая значимость и перспективы развития.

\textbf{В приложениях} могут быть приведены листинги конфигураций,
таблицы параметров прогонов (включая JSON\-метаданные экспериментов),
дополнительные графики и описание сценариев запуска.

\end{document}

% либо автономная сборка: pdflatex magister_tsel_zadachi_struktura.tex
%
\documentclass[12pt,a4paper]{article}
\usepackage[utf8]{inputenc}
\usepackage[T2A]{fontenc}
\usepackage[russian]{babel}
\usepackage{geometry}
\geometry{left=30mm,right=15mm,top=20mm,bottom=20mm}
\usepackage{indentfirst}
\setlength{\parindent}{1.25cm}

\begin{document}

\section*{Цель исследования}

\textbf{Целью} магистерской диссертации является разработка и развёртка
имитационной модели ансамбля беспилотных мобильных платформ,
позволяющей воспроизводить типовые режимы группового полёта и
анализировать устойчивость, координацию и качество управления при
изменении числа агентов, геометрии постановки задачи и параметров
законов управления и телеметрии (в том числе в условиях ограничений
канала связи и шумов измерений), с оформлением воспроизводимого
конвейера экспериментов и представлением результатов в виде
логов, метаданных прогона и графической визуализации.

\section*{Задачи исследования}

Для достижения указанной цели в работе ставятся следующие \textbf{задачи}:
\begin{enumerate}
  \item Провести обзор литературы и нормативно\-технической базы по
    моделированию многоагентных БПЛА, программно\-интегрированным
    стендам SITL/HIL, протоколу MAVLink и методам координации роя
    (лидер\,--\,последователь, формации, локальные законы соседства).
  \item Сформулировать постановку задачи имитации ансамбля: состав
    подсистем (динамика носителя в среде симуляции, обмен данными,
    сценарий верхнего уровня), требования к масштабируемости
    (число платформ, сетевые порты, начальные положения, длительность
    эксперимента, файлы параметров автопилота).
  \item Обосновать и описать используемые методы: представление
    состояния в локальной системе координат (NED), извлечение
    оценок положения из телеметрии (включая SIM\_STATE и связанные
    сообщения), преобразования «географические координаты\,--\,NED»,
    организация потокобезопасного доступа к соединению MAVLink,
    алгоритмы законов управления в сценариях (в частности,
    пропорционально\-интегральные и кусочно\-непрерывные/сигмоидные
    отображения ошибки в команды RC override), критерии и метрики
    анализа (расстояния между агентами, скорости, PWM, параметры
    устойчивости относительно опорных узлов).
  \item Реализовать программный комплекс имитации: запуск нескольких
    экземпляров SITL (ArduPilot), сценарии на языке Python,
    ядро взаимодействия с автопилотом по MAVLink, режимы с
    опциональной 3D\-визуализацией (Webots) и 2D\-визуализацией,
    пакетный запуск серии экспериментов с фиксацией метаданных.
  \item Выполнить серию вычислительных экспериментов по выбранным
    сценариям (включая сценарии сравнения при наличии шума в канале
    или в оценке положения), верифицировать корректность развёртки
    и воспроизводимость результатов.
  \item Систематизировать полученные результаты, оформить выводы
    о применимости разработанной имитационной модели для анализа
    поведения ансамбля и направлениях дальнейшего развития
    (расширение класса сценариев, учёт более сложных моделей среды
    и связи).
\end{enumerate}

\section*{Структура работы}

Диссертация состоит из введения, трёх глав, заключения, списка
использованных источников и приложений. Ниже приведена логическая
связка глав: в главе~2 излагаются методы и модели, в главе~3\,---
их программная реализация и результаты экспериментов по тем же
разделам (теория\,--\,практика).

\textbf{Во введении} обосновывается актуальность темы, формулируются
цель и задачи, кратко характеризуется объект и предмет исследования,
перечисляются научная новизна и практическая значимость, приводится
обзор структуры работы и положений, выносимых на защиту.

\textbf{Глава~1. Обзор литературы в рамках диссертации}
\begin{enumerate}
  \item Состояние исследований по координации ансамблей БПЛА:
    централизованные и децентрализованные подходы, устойчивость
    и ограничения по связи.
  \item Программно\-аппаратные стенды и имитация: SITL/HIL,
    интеграция с Gazebo/Webots, типовые стеки автопилота.
  \item Протокол MAVLink, модели телеметрии и команд верхнего
    уровня; вопросы синхронизации и частоты обновления данных.
  \item Методы анализа результатов имитации: метрики близости
    к заданной конфигурации, безопасные интервалы, сравнение
    вариантов алгоритма и параметров.
\end{enumerate}

\textbf{Глава~2. Обзор использованных методов, математики и алгоритмов}
\begin{enumerate}
  \item Математическое описание кинематики и используемых систем
    координат; переход от географических координат к локальному NED.
  \item Модель информационного обмена в имитационном стенде:
    потоки MAVLink, выбор сообщений для оценки состояния,
    организация параллельного доступа (потокобезопасный обработчик
    соединения).
  \item Класс рассматриваемых законов управления ансамблем:
    сценарии лидер\,--\,последователь, формации, цепочки с опорными
    узлами, локальные законы «антенна/соседство»; параметризация
    законов (коэффициенты ПИ, ограничения PWM, радиусы видимости,
    шаг дискретизации по PWM и~т.\,п.).
  \item Постановка вычислительного эксперимента: воспроизводимость,
    масштабирование числа агентов, пакетные прогоны, фиксация
    метаданных и логов для последующего анализа.
  \item Методы оценки качества: показатели по траекториям, отклонениям
    от якоря, согласованности PWM, сравнение режимов «без шума /
    со шумом» (при наличии соответствующего сценария).
\end{enumerate}

\textbf{Глава~3. Практическая реализация алгоритмов, фрагменты кода,
визуализация и результаты}
\begin{enumerate}
  \item Архитектура программного комплекса: сценарный уровень,
    модуль взаимодействия с MAVLink, запуск SITL и распределение
    сетевых портов; конфигурация параметров автопилота.
  \item Реализация геометрических и навигационных преобразований
    и потокобезопасного приёма/передачи сообщений (иллюстрация
    ключевыми фрагментами кода и пояснениями).
  \item Реализация выбранных сценариев ансамбля: соответствие
    алгоритмам главы~2, настраиваемые параметры сценария,
    журналирование телеметрии в файлы эксперимента.
  \item Средства визуализации (2D/3D) и конвейер постобработки:
    построение графиков траекторий, скоростей, PWM и других
    величин по сохранённым логам.
  \item Результаты вычислительных экспериментов и их интерпретация
    в контексте цели работы; обсуждение ограничений модели.
\end{enumerate}

\textbf{В заключении} формулируются основные результаты, степень
достижения цели, практическая значимость и перспективы развития.

\textbf{В приложениях} могут быть приведены листинги конфигураций,
таблицы параметров прогонов (включая JSON\-метаданные экспериментов),
дополнительные графики и описание сценариев запуска.

\end{document}

% либо автономная сборка: pdflatex magister_tsel_zadachi_struktura.tex
%
\documentclass[12pt,a4paper]{article}
\usepackage[utf8]{inputenc}
\usepackage[T2A]{fontenc}
\usepackage[russian]{babel}
\usepackage{geometry}
\geometry{left=30mm,right=15mm,top=20mm,bottom=20mm}
\usepackage{indentfirst}
\setlength{\parindent}{1.25cm}

\begin{document}

\section*{Цель исследования}

\textbf{Целью} магистерской диссертации является разработка и развёртка
имитационной модели ансамбля беспилотных мобильных платформ,
позволяющей воспроизводить типовые режимы группового полёта и
анализировать устойчивость, координацию и качество управления при
изменении числа агентов, геометрии постановки задачи и параметров
законов управления и телеметрии (в том числе в условиях ограничений
канала связи и шумов измерений), с оформлением воспроизводимого
конвейера экспериментов и представлением результатов в виде
логов, метаданных прогона и графической визуализации.

\section*{Задачи исследования}

Для достижения указанной цели в работе ставятся следующие \textbf{задачи}:
\begin{enumerate}
  \item Провести обзор литературы и нормативно\-технической базы по
    моделированию многоагентных БПЛА, программно\-интегрированным
    стендам SITL/HIL, протоколу MAVLink и методам координации роя
    (лидер\,--\,последователь, формации, локальные законы соседства).
  \item Сформулировать постановку задачи имитации ансамбля: состав
    подсистем (динамика носителя в среде симуляции, обмен данными,
    сценарий верхнего уровня), требования к масштабируемости
    (число платформ, сетевые порты, начальные положения, длительность
    эксперимента, файлы параметров автопилота).
  \item Обосновать и описать используемые методы: представление
    состояния в локальной системе координат (NED), извлечение
    оценок положения из телеметрии (включая SIM\_STATE и связанные
    сообщения), преобразования «географические координаты\,--\,NED»,
    организация потокобезопасного доступа к соединению MAVLink,
    алгоритмы законов управления в сценариях (в частности,
    пропорционально\-интегральные и кусочно\-непрерывные/сигмоидные
    отображения ошибки в команды RC override), критерии и метрики
    анализа (расстояния между агентами, скорости, PWM, параметры
    устойчивости относительно опорных узлов).
  \item Реализовать программный комплекс имитации: запуск нескольких
    экземпляров SITL (ArduPilot), сценарии на языке Python,
    ядро взаимодействия с автопилотом по MAVLink, режимы с
    опциональной 3D\-визуализацией (Webots) и 2D\-визуализацией,
    пакетный запуск серии экспериментов с фиксацией метаданных.
  \item Выполнить серию вычислительных экспериментов по выбранным
    сценариям (включая сценарии сравнения при наличии шума в канале
    или в оценке положения), верифицировать корректность развёртки
    и воспроизводимость результатов.
  \item Систематизировать полученные результаты, оформить выводы
    о применимости разработанной имитационной модели для анализа
    поведения ансамбля и направлениях дальнейшего развития
    (расширение класса сценариев, учёт более сложных моделей среды
    и связи).
\end{enumerate}

\section*{Структура работы}

Диссертация состоит из введения, трёх глав, заключения, списка
использованных источников и приложений. Ниже приведена логическая
связка глав: в главе~2 излагаются методы и модели, в главе~3\,---
их программная реализация и результаты экспериментов по тем же
разделам (теория\,--\,практика).

\textbf{Во введении} обосновывается актуальность темы, формулируются
цель и задачи, кратко характеризуется объект и предмет исследования,
перечисляются научная новизна и практическая значимость, приводится
обзор структуры работы и положений, выносимых на защиту.

\textbf{Глава~1. Обзор литературы в рамках диссертации}
\begin{enumerate}
  \item Состояние исследований по координации ансамблей БПЛА:
    централизованные и децентрализованные подходы, устойчивость
    и ограничения по связи.
  \item Программно\-аппаратные стенды и имитация: SITL/HIL,
    интеграция с Gazebo/Webots, типовые стеки автопилота.
  \item Протокол MAVLink, модели телеметрии и команд верхнего
    уровня; вопросы синхронизации и частоты обновления данных.
  \item Методы анализа результатов имитации: метрики близости
    к заданной конфигурации, безопасные интервалы, сравнение
    вариантов алгоритма и параметров.
\end{enumerate}

\textbf{Глава~2. Обзор использованных методов, математики и алгоритмов}
\begin{enumerate}
  \item Математическое описание кинематики и используемых систем
    координат; переход от географических координат к локальному NED.
  \item Модель информационного обмена в имитационном стенде:
    потоки MAVLink, выбор сообщений для оценки состояния,
    организация параллельного доступа (потокобезопасный обработчик
    соединения).
  \item Класс рассматриваемых законов управления ансамблем:
    сценарии лидер\,--\,последователь, формации, цепочки с опорными
    узлами, локальные законы «антенна/соседство»; параметризация
    законов (коэффициенты ПИ, ограничения PWM, радиусы видимости,
    шаг дискретизации по PWM и~т.\,п.).
  \item Постановка вычислительного эксперимента: воспроизводимость,
    масштабирование числа агентов, пакетные прогоны, фиксация
    метаданных и логов для последующего анализа.
  \item Методы оценки качества: показатели по траекториям, отклонениям
    от якоря, согласованности PWM, сравнение режимов «без шума /
    со шумом» (при наличии соответствующего сценария).
\end{enumerate}

\textbf{Глава~3. Практическая реализация алгоритмов, фрагменты кода,
визуализация и результаты}
\begin{enumerate}
  \item Архитектура программного комплекса: сценарный уровень,
    модуль взаимодействия с MAVLink, запуск SITL и распределение
    сетевых портов; конфигурация параметров автопилота.
  \item Реализация геометрических и навигационных преобразований
    и потокобезопасного приёма/передачи сообщений (иллюстрация
    ключевыми фрагментами кода и пояснениями).
  \item Реализация выбранных сценариев ансамбля: соответствие
    алгоритмам главы~2, настраиваемые параметры сценария,
    журналирование телеметрии в файлы эксперимента.
  \item Средства визуализации (2D/3D) и конвейер постобработки:
    построение графиков траекторий, скоростей, PWM и других
    величин по сохранённым логам.
  \item Результаты вычислительных экспериментов и их интерпретация
    в контексте цели работы; обсуждение ограничений модели.
\end{enumerate}

\textbf{В заключении} формулируются основные результаты, степень
достижения цели, практическая значимость и перспективы развития.

\textbf{В приложениях} могут быть приведены листинги конфигураций,
таблицы параметров прогонов (включая JSON\-метаданные экспериментов),
дополнительные графики и описание сценариев запуска.

\end{document}
